% ===========================================================
% discussion.tex
% ===========================================================
\section{Discussion}
\label{sec:discussion}

The pipeline's strengths are also its limitations: dependence on
public CEN, SCVIC, and PLEXOS feeds confers auditability and
reproducibility but also exposes it to publishing-agency cadence,
completeness, and rate limits.

\paragraph{Rate-limit dependency.}
The CEN SIP API enforces a soft \SI{8}{requests\per\second} limit on
the anonymous-read key and intermittently returns HTTP~502 on heavy
endpoints (\texttt{cmg-programado-pcp}, per-bar\_transf
\texttt{cmg-real} fan-out).  The pipeline mitigates with parallel
fetching, bounded concurrency, and an aggressive parquet cache; a
cold-cache day requires \SI{75}{\second} of wall time and an active
session.

\paragraph{Hand-curated bus aliases.}
The bus-name bridge of Section~\ref{sec:inputs} relies on \num{29}
hand-curated overrides to reconcile PLEXOS node names that the
heuristic alias generator cannot match unambiguously.  Maintenance
burden grows linearly with the catalogue.

\paragraph{SCVIC tier wedge.}
The \texttt{tier\_uncertainty} flag identifies cells where the picker
selects a CV tier whose next-up neighbour is more than \$2/MWh away.
These cells (\SI{13.7}{\percent} of the 2026-04-22 sample) carry
most of the worst-case error.  Mitigation requires either a finer
SCVIC ladder (out of pipeline control) or a within-tier interpolation
rule, the latter blocked by the absence of access to the
PLEXOS-internal piecewise-linear cost segment binding at LP solve.

\paragraph{Loss-factor decomposition.}
The per-bus loss correction is not a proper PJM-style
energy/loss/congestion decomposition: CEN does not publish the
per-unit Power Transfer Distribution Factor matrix.  The PID-extracted
MLF approximates the loss component but does not isolate the
congestion component; the pipeline therefore cannot answer
counterfactual questions such as ``how would $\hat\lambda$ at bus
$b$ change if line $\ell$ were de-rated?''~— a question for which
\texttt{gtopt} SDDP is well-equipped (Section~\ref{sec:conclusions}).

\paragraph{Reservoir-hydro as marginal.}
The pipeline currently excludes reservoir-hydro plants from the
candidate pool because the published feeds do not expose the
implicit water value at which the LP cleared.  The natural mitigation
is to seed reservoir-hydros with the
\texttt{cmg-programado-pcp} median as a synthetic CV; the pipeline
is structurally ready, blocked on stable availability of that
endpoint.

\paragraph{Single-pollutant scope.}
The current emission factor is CO$_2$ only.  Multi-pollutant
extension (NO$_x$, SO$_2$, PM$_{2.5}$) requires per-unit declared
factors from the Registro de Emisiones y Transferencias de
Contaminantes (RETC); the integration is straightforward
(\texttt{FUEL\_EF} becomes \texttt{POLLUTANT\_EF[fuel][pollutant]})
and would close the last methodological gap relative to PJM~LME.

\paragraph{Carbon-credit implications.}
A per-bus per-quarter $\varepsilon^{\mathrm{bus}}_{b,t}$ is the
formal input expected by CDM Tool~07 OM-DDA~\cite{UNFCCCTool07},
Verra VCS v4.5~\cite{VerraVCS45}, EnergyTag granular
certificates~\cite{EnergyTagGC}, the 24/7~CFE
Compact~\cite{CFE247}, and the Article~6.4 SBM 4DM
draft~\cite{Article64Draft}; the 15-minute output natively
satisfies all of their granularity requirements.  A worked example
illustrates the magnitude of the methodology choice: a hypothetical
\SI{100}{\MW} PV plant at CRUCERO 220~kV on 2026-04-22 yields
\SI{19.9}{\tonne CO_2}/day under hourly-marginal accounting,
versus \SI{270.6}{\tonne CO_2} under OM-DDA and \SI{258.4}{\tonne CO_2}
under the national-average factor~\cite{InvNac2024} — a thirteen-fold
gap, because the marginal unit at CRUCERO during the
\SIrange{09}{16}{\hour} window is itself renewable
($\varepsilon^{\mathrm{bus}}=0$) and a new PV injection at those
hours displaces zero-emission energy.  The same reconstruction
underwrites the natural BESS-arbitrage extension: a battery that
charges in the low-$\varepsilon$ window and discharges at the
high-$\varepsilon$ evening peak earns nearly the full
$\Delta\varepsilon$ as ER per discharged MWh.
